% Осень26!A25; нумерация внутри столбца без пустых строк.
% Источник: лекция 3, §9 (строка 251)
% Источник: лекция 3, §10 (строка 272)
% Источник: лекция 3, §11 (строка 298)
\ticket{22}{Лемма о структуре открытых множеств и множеств меры 0}
\label{ticket:22}

\subsection*{Лебеговская σ-алгебра и топология; лемма о структуре открытых множеств и множествах меры 0 --- формулировки}

\textbf{Свойство 2.} Лебеговская σ-алгебра \(\mathcal{M}^m\) «дружит с топологией»: содержит \textbf{все открытые} и \textbf{все замкнутые} множества \(\mathbb{R}^m\).

Замкнутые --- дополнения открытых, а σ-алгебра содержит дополнения; достаточно проверить открытые. А открытые содержатся, «потому что открытое множество можно разбить на ячейки» --- короткий ответ; технически это лемма:

\begin{lemma}

\emph{1 (о структуре открытых множеств.} Пусть \(O\subset\mathbb{R}^m\) --- открытое множество. Тогда \[O=\bigsqcup_i Q_i\] --- счётное (не более чем счётное) \textbf{дизъюнктное} объединение ячеек, причём дополнительно можно потребовать (каждую из опций или только некоторые --- «как хотите»):

\begin{itemize}
\tightlist
\item
  ячейки \(Q_i\) --- \textbf{кубы} (все рёбра равны);
\item
  координаты вершин \(Q_i\) \textbf{рациональны} (все \(a_i,b_i\in\mathbb{Q}\); более сильная опция --- \textbf{двоично-рациональны}: вида \(k/2^s\));
\item
  ячейки не соприкасаются с границей: \(\bar Q_i\subset O\) (замыкание ячейки лежит в \(O\); ячейки «запрятаны внутрь», у границы измельчаются, но до границы не дотягиваются).
\end{itemize}

\end{lemma}

\begin{lemma}

\emph{2 (о множествах меры 0.} Пусть \(E\in\mathcal{M}^m\). Тогда \(\lambda E=0\) в том и только в том случае, если \[\forall\varepsilon>0\ \exists\ \text{ячейки } P_k:\qquad E\subset\bigcup_k P_k,\qquad \sum_k\lambda P_k<\varepsilon .\] Вариант: вместо ячеек можно брать \textbf{шары} \(B_i\): \(E\subset\bigcup_iB_i\), \(\sum_i\lambda B_i<\varepsilon\cdot m^{m/2}\).

«Любое множество меры ноль можно накрыть чрезвычайно маленькими клеточками, суммарная мера которых очень маленькая --- и вы закрыли множество меры ноль».

\end{lemma}

\subsection*{Доказательство леммы о структуре открытых множеств}

\textbf{Требует сверки.} Между формулировками и доказательством в записи --- перерыв (содержательного материала там нет).

\textbf{Шаг 1 (индивидуальные ячейки).} Для каждой точки \(x\in O\) в силу открытости найдётся её «личная индивидуальная» ячейка \(Q_x\): кубическая, с двоично-рациональными координатами, такая что \[x\in Q_x,\qquad \bar Q_x\subset O .\] Тогда очевидно \[O=\bigcup_{x\in O}Q_x\] (левая часть содержится в правой: каждая точка лежит хотя бы в своей ячейке; правая в левой: каждая \(Q_x\subset O\)).

\textbf{Шаг 2 (прореживание до счётного).} Получилось объединение по континуальному множеству точек, но различных значений \(Q_x\) --- счётное число: рациональных ячеек не больше чем \(|\mathbb{Q}^{2m}|=\aleph_0\) (у ячейки два конца, каждый --- набор из \(m\) рациональных координат). «Содержательно в этом объединении присутствует лишь счётное множество ячеек, остальные идут с повторениями; проредим множество точек области --- оставим счётное количество представителей \(x_i\)» (теоретико-множественный фокус; «какая-нибудь аксиома выбора опять сыграет»): \[O=\bigcup_iQ_{x_i}.\] Отсюда, в частности, \textbf{все открытые множества измеримы}: счётное объединение ячеек (элементов полукольца), а σ-алгебра замкнута относительно счётных объединений.

\textbf{Шаг 3 (дизъюнктивизация).} Есть счётное объединение, а нужно \textbf{дизъюнктное}. Перенумеруем: \(Q_1=Q_{x_1}\). Добавляем \(Q_{x_2}\) --- но не целиком, а только новые точки: разность \(Q_{x_2}\setminus Q_1\). Дальше --- \(Q_{x_3}\) минус уже набранное, и т.д.:

\[O=Q_1\sqcup\bigl(Q_{x_2}\setminus Q_1\bigr)\sqcup\Bigl(Q_{x_3}\setminus\bigl(Q_1\cup(Q_{x_2}\setminus Q_1)\bigr)\Bigr)\sqcup\cdots\]

Ключевое свойство полукольца: \textbf{разность ячейки и конечного объединения ячеек --- конечное дизъюнктное объединение ячеек}. Но это дало бы просто ячейки, а мы хотим сохранить \textbf{кубичность}. Вот зачем двоично-рациональные координаты: все управляющие параметры картинки (координаты граней старого куба и «откусываемой» области) --- двоично-рациональные числа; приведём их к общему знаменателю \(2^s\) («какой-то из знаменателей \(2^{\text{степень}}\) самый крупный --- его и возьмём; можно подобрать \(s\), обслуживающее и все проекции на другие оси») и наложим «миллиметровочку» с клеткой со стороной \(1/2^s\): все вспомогательные линии лягут на сетку, и разность разрежется на конечное число \textbf{кубических} ячеек со стороной \(1/2^s\) (каждая клеточка квадратная: и горизонтальная, и вертикальная сторона \(1/2^s\)). Для \(m\)-мерной картинки обобщение «абсолютно один в один, только обозначения более громоздкие».

Индукция по шагам: \(Q_{x_k}\setminus\bigl(\bigcup_{j<k}\text{набранного}\bigr)\) --- конечное дизъюнктное объединение кубических ячеек, дизъюнктных со всем предыдущим. «Номера быстрорастут, но у нас бесконечно много номеров, не жалко; с каждым \(Q_i\) разберёмся. Счётное объединение победим». Итог: \(O=\bigsqcup_iQ_i\) --- счётное дизъюнктное объединение кубов.

\textbf{Проверка опций.} Ячейки остались кубическими и двоично-рациональными; каждое новое \(Q_i\) содержится в каком-то \(Q_{x_i}\) вместе с замыканием, а \(\bar Q_{x_i}\subset O\), поэтому и \(\bar Q_i\subset O\) --- «всё выполнилось». \(\blacksquare\)

\subsection*{Множества меры 0: доказательство; покрытие шарами}

\textbf{Пункт 2 леммы --- «очевиден по модулю основной теоремы»}, а именно п. 5 теоремы о продолжении: мера множества по определению считается как инфимум сумм по покрытиям ячейками: \[\lambda E=\inf\Bigl\{\sum_k\lambda P_k:\ E\subset\bigcup_kP_k,\ P_k\ \text{ячейки}\Bigr\}.\] Если \(\lambda E=0\), то инфимум нулевой, значит для любого \(\varepsilon>0\) найдётся покрытие с \(\sum_k\lambda P_k<\varepsilon\). Обратно: если для любого \(\varepsilon\) есть покрытие с суммой \(<\varepsilon\), то (полуаддитивность) \(\lambda E\le\sum_k\lambda P_k<\varepsilon\) для всех \(\varepsilon\), т.е. \(\lambda E=0\).

\textbf{Приведение к кубическим ячейкам.} Дана произвольная ячейка \(P_k\); заменим её конечным набором кубов, почти не увеличив меру: чуть-чуть \textbf{раздуем} \(P_k\) до двоично-рациональной ячейки \(\widetilde P_k\supset P_k\) («ту же биосекцию делали в доказательстве счётной аддитивности классического объёма»), контролируя \(\lambda(\widetilde P_k\setminus P_k)<\varepsilon/2^{k+1}\); затем разрежем \(\widetilde P_k\) по сетке с подходящим шагом \(1/2^s\) на кубики. Покрытие сохраняется, суммарная мера вырастает не более чем на \(\sum_k\varepsilon/2^{k+1}=\varepsilon/2\).

\textbf{Замена кубов шарами.} Пусть \(E\) покрыто кубами \(Q_i\) с \(\sum_i\lambda Q_i<\varepsilon\), где \(Q_i\) --- куб с центром \(a_i\) и полустороной \(r_i\) (сторона \(2r_i\), объём \((2r_i)^m\)). Заменяем каждый куб \textbf{описанным шаром}: \[Q_i\subset B(a_i,r_i\sqrt m)\] (теорема Пифагора: \(m\) ортогональных катетов длины \(r_i\) каждый). «Меру шара я, к сожалению, не знаю --- но она мне и не нужна: шар содержится в кубе с тем же центром и полустороной \(r_i\sqrt m\), а его объём я знаю». По монотонности меры: \[\lambda B(a_i,r_i\sqrt m)\le(2r_i\sqrt m)^m=(2r_i)^m\,m^{m/2}=\lambda Q_i\cdot m^{m/2}.\] Суммируя: \(E\subset\bigcup_iB_i\) и \[\sum_i\lambda B_i\le m^{m/2}\sum_i\lambda Q_i<\varepsilon\, m^{m/2}.\] «Если хотите покрытие шарами с суммарной мерой меньше тысячной --- постройте исходное кубическое покрытие с \(\varepsilon/1000\), делённым на \(m^{m/2}\); шары раздуются --- всё сойдётся». \(\blacksquare\)
